problem:  
Let \((M^{2n}, g, J)\) be a compact Kähler manifold with **positive orthogonal Ricci curvature** (that is, for every nonzero \((1,0)\)-vector \(X\),  
\[
\mathrm{Ric}^\perp(X,\overline{X}) := \mathrm{Ric}(X,\overline{X}) - \frac{R(X,\overline{X},X,\overline{X})}{|X|^2} > 0.
\]  
)  
It is known (Yang ’19, Chen–Yang–Zheng ’21) that under this curvature condition:
- any two distinct **divisors** must intersect; and  
- several Frankel-type intersection results hold in low dimensions.  

However, no general result is known for higher codimension.

**Open problem:**  
Does positive orthogonal Ricci curvature necessarily enforce the **Frankel intersection property** for all pairs of closed complex submanifolds \(X^p, Y^q \subset M^{2n}\) with \(p+q\ge n\)?  

If not, determine the optimal dimension threshold function \(f(n)\) such that  
\[
p+q \ge f(n) \implies X \cap Y \neq \varnothing
\]  
for all such submanifolds in Kähler manifolds with \(\mathrm{Ric}^\perp>0\).  

Is \(f(n)=n\) universally valid, or does the critical threshold depend on the finer structure of the curvature tensor—e.g., whether \(\operatorname{Rm}(M)\) lies in a given \(U(n)\)-invariant curvature cone between those for holomorphic bisectional and orthogonal Ricci positivity?

Why is it a "good" problem:  
This problem sits precisely at the **boundary between analytic curvature inequalities and global intersection rigidity**. While Frankel’s theorem under positive holomorphic bisectional curvature fully enforces intersection, and the divisor case under \(\mathrm{Ric}^\perp>0\) is known, the higher-codimension situation is **wide open**. The orthogonal Ricci condition captures *average holomorphic curvature positivity*—a subtle weakening that retains strong metric implications yet may no longer control geodesic convexity in all complex directions. Understanding whether it still compels intersection phenomena would clarify the minimal analytic positivity required for topological rigidity.  

The introduction of a threshold \(f(n)\) broadens this to a quantitative curvature–topology question: can global intersection rigidity decay continuously with curvature weakening? Resolving this could reveal a previously unnoticed **spectrum of Frankel-type phenomena** governed by representation-theoretic curvature cones. Conceptually simple but deeply connected to several modern geometric themes—curvature cone analysis, Kähler rigidity, and global geodesic geometry—this problem exemplifies the combination of simplicity, depth, and cross-field implications that define a “good” research problem.